\subsection{Sparsification}
\label{sec:method:reduction:sparsification}
Two of the four methods act on the edge set alone ($\eta_{\mathcal{V}} = 1$). The other two
select a node subset $\mathcal{V}' \subseteq \mathcal{V}$ and return the induced subgraph
$G[\mathcal{V}']$, removing incident edges too (\ref{sec:prelim:reduction:measuring}). Only
one of the four has a compression parameter. For the rest, retention is a measurement,
reported in~\ref{sec:results:rq2:offline}.
% Original: Two of the four methods below act on the edge set alone, replacing $\mathcal{E}$ with a subset
% $\mathcal{E}' \subseteq \mathcal{E}$ and leaving $\eta_{\mathcal{V}} = 1$. The other two
% select a node subset $\mathcal{V}' \subseteq \mathcal{V}$ and return the induced subgraph
% $R(G) = G[\mathcal{V}']$, so every edge with an endpoint outside $\mathcal{V}'$ is removed as
% well. That difference determines which retention figure is meaningful for each method, and it
% is why~\ref{sec:prelim:reduction:measuring} insists on reporting node and edge retention
% separately rather than as one compression number. Only one of the four has a compression
% parameter at all; for the other three the achieved retention is a measurement, reported
% in~\ref{sec:results:rq2:offline} rather than here.

\subsubsection{Random Edge Dropout}
\label{sec:method:sparse:dropout}
This naive baseline uniformly drops a random fraction of edges, a control for structural
degradation (\ref{sec:relwork:reduction:sparsification}): each edge survives when an
independent uniform draw clears the drop rate $p$,
\begin{equation}
    \mathcal{E}' = \{\, e \in \mathcal{E} : U_e \geq p \,\},
    \qquad
    U_e \overset{\text{i.i.d.}}{\sim} \mathrm{Unif}[0,1],
    \label{eq:method:sparse-dropout}
\end{equation}
so $\mathbb{E}[\eta_{\mathcal{E}}] = 1 - p$, with the mask drawn once and fixed rather than
redrawn per epoch \citep{rong2020dropedge}.
% Original: This naive baseline uniformly drops a fixed random fraction of edges across the graph,
% acting as a control for structural degradation. Each edge draws an independent uniform
% variate and survives when the draw clears the drop rate $p$,
% [[equation unchanged, see above]]
% so each edge is kept independently with probability $1 - p$ and
% $\mathbb{E}[\eta_{\mathcal{E}}] = 1 - p$. Random edge removal is established as a per-epoch
% regulariser for deep GNN training \citep{rong2020dropedge}. Here the
% mask of~\eqref{eq:method:sparse-dropout} is instead drawn once and fixed, which makes it a
% reduction rather than a regulariser.

The drop rate is the family's only tunable parameter, so its value decides which
parameter-free method it matches, and it was set for a pairing other than the one it is used
in. It is fixed at $p = 0.3$, chosen to place $\mathbb{E}[\eta_{\mathcal{E}}]$ near the
AND-gate method's edge retention (\ref{sec:method:sparse:andgate}). The pairing it actually
serves is against the spanning forest (\ref{sec:method:sparse:forest}), which requires
$p = 1 - \eta_{\mathcal{E}}^{\mathrm{forest}} \approx 0.419$ instead: at the reported $0.3$,
that pairing removes visibly less than its partner, so no accuracy difference inside it can
be read as a difference between the two methods (\ref{sec:results:rq4:pairings},
\ref{fig:rq4_pairing_compression}).
% TODO: rerun random edge dropout at 0.419, or drop the spanning-forest pairing. (I think i will jus add this as an eval like oh that could have been better. )
% Original: The drop rate is the only freely tunable compression parameter in the family, so its value
% decides which parameter-free method it can be matched against, and it was set for a pairing
% other than the one it is used in. It is fixed at $p = 0.3$, chosen to place
% $\mathbb{E}[\eta_{\mathcal{E}}]$ near the edge retention of the AND-gate method
% of~\ref{sec:method:sparse:andgate}. The pairing it actually serves is against the spanning
% forest of~\ref{sec:method:sparse:forest}, which requires
% $p = 1 - \eta_{\mathcal{E}}^{\mathrm{forest}} \approx 0.419$ instead. The reported runs use
% $0.3$, so that pairing removes visibly less than its partner and no accuracy difference
% inside it can be read as a difference between the two methods
% (\ref{sec:results:rq4:pairings}, \ref{fig:rq4_pairing_compression}).

\subsubsection{AND-Gate Retention}
\label{sec:method:sparse:andgate}
This domain-specific method removes the circuit interface and keeps the logic: in the node
role notation of~\eqref{eq:prelim:role},
\begin{equation}
    \mathcal{V}' = \{\, v \in \mathcal{V} :
        \tau(v) \notin \{\mathrm{PI}, \mathrm{PO}\} \,\},
    \label{eq:method:sparse-andgate}
\end{equation}
so all PI and PO nodes are dropped with every incident edge, and the internal AND network
survives intact. The rationale: the PI/PO boundary is fixed by the circuit's interface and
cannot be restructured by synthesis, so the AND network carries the optimizable structure.
The method is parameter-free. Its node retention is exactly the AND-gate share of the
circuit, a corpus statistic (\ref{sec:method:data:statistics:scale}) rather than a dial, and
it varies from graph to graph.
% Original: This domain-specific method removes the circuit interface and keeps the logic. In the node
% role notation of~\eqref{eq:prelim:role},
% [[equation unchanged, see above]]
% so all PI and PO nodes are dropped with every edge that
% touches them, and the internal AND network survives intact. The rationale is that the PI/PO
% boundary is fixed by the circuit's interface and cannot be restructured by synthesis, so the
% AND network carries the optimizable structure. The method is parameter-free. Its node
% retention is exactly the AND-gate share of the circuit, a corpus statistic
% (\ref{sec:method:data:statistics:scale}) rather than a value that can be dialled, and it
% varies from graph to graph.

\subsubsection{Random Spanning Forest}
\label{sec:method:sparse:forest}
This method extracts a sparse connectivity backbone
(\ref{sec:relwork:reduction:sparsification}). Let $\tilde{G}$ be the undirected projection
of the AIG, with each edge an independent weight $w_e \sim \mathrm{Unif}[0,1]$. The method
keeps the minimum-weight spanning forest
\begin{equation}
    F = \operatorname*{arg\,min}_{F' \in \mathfrak{F}(\tilde{G})}
        \sum_{e \in F'} w_e,
    \qquad
    |\mathcal{E}(F)| = |\mathcal{V}| - c,
    \label{eq:method:sparse-forest}
\end{equation}
with $\mathfrak{F}(\tilde{G})$ the set of spanning forests of $\tilde{G}$ and $c$ the number
of connected components, computed with Kruskal's algorithm \citep{kruskal1956mst}. A forest
contains no cycles, so every reconvergent path loses an edge while connectivity stays
intact. Because $F$ is computed on $\tilde{G}$, edge direction plays no part in which edges
survive, a real loss of information on a DAG and one reason to expect this method to damage
the predictive signal.
% Original: This method extracts a sparse connectivity backbone. Let $\tilde{G}$ be the undirected
% projection of the AIG, and assign each edge an independent weight
% $w_e \sim \mathrm{Unif}[0,1]$. The method keeps the minimum-weight spanning forest
% [[equation unchanged, see above]]
% where $\mathfrak{F}(\tilde{G})$ is the set of spanning forests of $\tilde{G}$ and $c$ the
% number of connected components, computed with Kruskal's algorithm \citep{kruskal1956mst}. The
% weights are almost surely distinct, so each draw has a unique minimiser and the randomised
% weights select a random forest rather than a deterministic one. A forest contains no cycles.
% Every reconvergent path therefore loses an edge, while connectivity is kept intact. Because
% $F$ is computed on $\tilde{G}$, edge direction plays no part in which edges survive. That is
% a real loss of information on a DAG, and one reason to expect this method to damage the
% predictive signal.

The edge count in~\eqref{eq:method:sparse-forest} also makes the method parameter-free:
$|\mathcal{V}| - c$ edges is a property of the graph rather than a setting, so its retention
follows from the density of the corpus (\ref{sec:results:rq2:offline}).

This method replaced a stretch-bounded spanner \citep{peleg1989spanners}
(\ref{sec:relwork:reduction:sparsification}), implemented and abandoned: AIGs have little of
the dense short-range redundancy a $t$-spanner exploits, so at stretch $t = 3$ the
randomised construction of \citet{baswana2007spanners} returned the graph nearly unchanged,
the smallest stretch the construction admits above the exact case $t = 1$, so a weaker
setting could not have avoided the failure. The result is worth reporting: a small, concrete
instance of the thesis's general claim that generic graph machinery does not transfer to
AIGs unexamined.
% Original: This method replaced a stretch-bounded spanner \citep{peleg1989spanners}, which was
% implemented and abandoned. A $t$-spanner keeps a subgraph $H \subseteq \tilde{G}$ with
% $d_H(u,v) \leq t \cdot d_{\tilde{G}}(u,v)$ for every node pair, so an edge can be discarded
% only when an alternative path of length at most $t$ survives. Spanners exploit dense
% short-range redundancy, and AIGs have little of it: the reconvergence they are rich in runs
% through long paths, so at stretch $t = 3$ the randomised construction of
% \citet{baswana2007spanners} returned the graph nearly unchanged. That stretch is the smallest
% the construction admits above the exact case $t = 1$, so the failure is not one a weaker
% setting could have avoided. The negative result is worth reporting: it is a small, concrete
% instance of the thesis's general claim that generic graph machinery does not transfer to AIGs
% unexamined.

\subsubsection{PageRank Node Pruning}
\label{sec:method:sparse:pagerank}
This centrality-based approach scores nodes by PageRank and retains the highest-scoring
fraction, dropping every incident edge too (\ref{sec:relwork:reduction:sparsification}). The
score is the stationary solution of
\begin{equation}
    \operatorname{pr}(v) = \frac{1 - \alpha}{|\mathcal{V}|}
    + \alpha \sum_{u \in \operatorname{fi}(v)}
        \frac{\operatorname{pr}(u)}{|\operatorname{fo}(u)|},
    \label{eq:method:sparse-pagerank}
\end{equation}
computed with fanin and fanout as in~\eqref{eq:prelim:fanin}, the mass of the POs (whose
fanout is empty) redistributed uniformly \citep{brin1998anatomy,langville2003pagerank},
keeping the $\lceil \kappa \, |\mathcal{V}| \rceil$ highest-scoring nodes. The damping factor
is left at the original $\alpha = 0.85$ \citep{brin1998anatomy}, untuned since it governs how
far centrality propagates, not how much the method removes.
% Original: This centrality-based approach scores every node by PageRank and retains only the
% highest-scoring fraction, implicitly dropping all edges incident to the removed
% low-centrality nodes. The score is the stationary solution of
% [[equation unchanged, see above]]
% computed on the directed graph with fanin and fanout as in~\eqref{eq:prelim:fanin} and the
% mass of the POs, whose fanout is empty, redistributed uniformly
% \citep{brin1998anatomy,langville2003pagerank}. The kept set $\mathcal{V}'$ collects the
% $\lceil \kappa \, |\mathcal{V}| \rceil$ nodes of highest score. The damping factor is left at
% the $\alpha = 0.85$ of the original formulation \citep{brin1998anatomy} and was not tuned,
% since it governs how far centrality propagates and not how much the method removes.

The retention fraction $\kappa$ is the method's compression knob, set to $\kappa = 0.8$,
calibrated to the AND-gate share (\ref{sec:method:data:statistics:scale}) so the two form a
generic and domain-informed pair at approximately equal node retention
(\ref{sec:results:rq4:pairings}). The equality is weaker than it looks: $\kappa$ fixes node
retention exactly on every graph, whereas the AND-gate share varies by circuit, so the pair
meets at the corpus mean, not graph by graph.
% Original: The retention fraction $\kappa$ is the method's compression knob and is set to $\kappa = 0.8$,
% calibrated to the AND-gate share of~\ref{sec:method:data:statistics:scale} so that the two
% form a generic and domain-informed pair at approximately equal node retention
% (\ref{sec:results:rq4:pairings}). The equality is weaker than the shared figure suggests.
% Here $\kappa$ fixes node retention exactly on every graph, whereas the AND-gate share varies
% by circuit, so the pair meets at the corpus mean and not graph by graph.

%TODO make sure these are very formulaic an dmathematical 